\heading{Local and domain-global probabilities}
At a mass fixed independently of the data, the local probability is
\begin{equation}
 p_{\rm loc}(m)=\Pr_{0}\{T(m)\geq T_{\rm obs}(m)\},
 \qquad Z_{\rm loc}(m)=\Phi^{-1}[1-p_{\rm loc}(m)],
\end{equation}
where larger $T$ is more signal-like and the probability includes equality.
This inclusive convention matters when a nonnegative fitted amplitude gives
a point mass at $T=0$: its inclusive tail is one. A displayed nonnegative
$Z$ is a plotting convention; the probability remains the primary quantity.
The common-mass alternative has one amplitude for each active dataset, so
$\sqrt{q_R}$ is not generally its local significance. Its null distribution
must be calibrated with the appropriate boundary and active dataset count.

For a declared mass grid $\mathcal G$ and a declared local ordering $S(m)$,
the scan statistic and its domain-global probability are
\begin{equation}
 M=\max_{m\in\mathcal G} S(m),\qquad
 p_{\mathcal G}=\Pr_0\{M\geq M_{\rm obs}\}.
\end{equation}
For example, $S=-\log p_{\rm loc}$ compares locally calibrated evidence
across masses; $S=q_R$ defines a different valid scan test. The observed
maximum and every null maximum must use the same ordering. A full-spectrum
toy supplies the entire correlated curve, rather than independent random
fluctuations at each mass. This accounts for the look-elsewhere effect within
that grid. See
\href{https://arxiv.org/abs/1005.1891}{Gross and Vitells} for scan inference and
\href{https://arxiv.org/abs/2206.12328}{Ananiev and Read} for correlated Gaussian
significance-field approximations.

Here the inherited full combined grid is $19,20,\ldots,250$ MeV. The active
supports are 19--100 MeV (2015), 39--180 MeV (2016), and 50--250 MeV
(2021, 10\%); only 50--100 MeV contains all three datasets. ``Common mass''
ties the mass of the active components and does not extend an experiment's
support. Adding sites to a domain cannot reduce its null exceedance
probability at a \emph{fixed threshold}. Comparisons of observed maxima can
also change the threshold. The number of grid points is not the number of
independent trials: nearby masses share events, background fits, and signal
templates. Finer sampling resolves maxima until grid convergence; a
1-MeV-grid calibration is not, by itself, a continuous-mass calibration.

The 80--105 MeV interval was chosen after inspection and is used only to
describe the reported structure. Its 2015 interval ends at 100 MeV.
Calibrating a statistic restricted to this interval answers a conditional
diagnostic question; it does not replace the full-domain discovery test or
correct the choice of this interval. Likewise, four width-specific global
probabilities are four conditional results, not a selection-corrected claim
based on the most favorable width. Such a claim would require repeating the
entire width/method/domain selection in joint null realizations. Reusing
the same fluctuations across widths preserves their strong dependence.

All probabilities here condition on the pinned observed-data-derived GP
source, its kernel policy, and the stated likelihood construction. Direct
Poisson scans repeat the local fits; Gaussian response fields approximate
their fluctuations. Neither propagates uncertainty in choosing or estimating
the source itself. Exact likelihood factorization therefore does not imply
unconditional physical-discovery calibration. Boundary asymptotics are
reviewed by \href{https://arxiv.org/abs/1007.1727}{Cowan et al.}; sparse toy
tails here are reported with exceedance counts and binomial bounds.
